\section{A distortion--dilation comparison}\label{sec:rate44}
The lower and upper profiles defined in Section~\ref{sec:profiles} do not
have identical quantifiers. They nevertheless obey a universal
inequality. Define the optimized per-command distortion rate
\begin{equation}\label{eq:rate44}
 \mathfrak r_{\calA}(k)=\sup_{B\ge1}
 \frac{-\log(1-\Gamma_{\calA}(k,B))}{B}.
\end{equation}
All logarithms in this paper are natural. The value $+\infty$ has its
usual extended meaning.

\begin{theorem}[Distortion bounds every common enclosure]\label{thm:rate44}
For every finite orthogonal alphabet and every feasible $k,a$ in
\eqref{eq:enclosure-profile},
\begin{equation}\label{eq:rate-inequality44}
 \mathfrak r_{\calA}(k)\le\log\lambda_{\calA}(k,a).
\end{equation}
In particular, every finite-word certificate gives the lower enclosure
bound $\lambda_{\calA}(k,a)\ge
(1-\Gamma_{\calA}(k,B))^{-1/B}$. The assertion does not require
commuting commands, a spectral gap, or any Diophantine hypothesis.
\end{theorem}
\begin{proof}
Take any feasible polytope $P=\conv(v_1,\ldots,v_k)$, contained in the
unit ball and containing $a$ times that ball, with common dilation
$\lambda$. Use the stochastic rows representing
$\lambda^{-1}U_bv_s$ by the same vertices. Fix a unit seed $u$ and
initialize the mean at $au$. This construction makes sense at every
length, without imposing a terminal decoder. For each complete word
$w$ of length $n$ it satisfies
\[
 \E[v_{S_n}\mid w]=a\lambda^{-n}U_wu.
\]
Let $Y_n=U_wu$ and $q_n=\E\|\E[Y_n\mid S_n]\|$. Pairing with the
bounded vector $v_{S_n}$ gives, under every word law,
\begin{equation}\label{eq:rate-lower-amplitude44}
 q_n\ge \E\ip{Y_n}{v_{S_n}}=a\lambda^{-n}.
\end{equation}
Fix $B$ and independently repeat an optimizing length-$B$ word law
$m$ times. The packet lemma applies at the $m$ endpoints, each of
which has at most $k$ labels. Since $q_0=1$,
\[
 a\lambda^{-mB}\le q_{mB}\le(1-\Gamma_{\calA}(k,B))^m.
\]
Taking $m$th roots and letting $m\to\infty$ removes the inradius
factor. Hence $\lambda^{-B}\le1-\Gamma_{\calA}(k,B)$. Take the
supremum over $B$ and then the minimum over enclosures. In particular,
$\Gamma=1$ is incompatible with a finite feasible enclosure at this
$k$, so no undefined logarithmic operation was used.
\end{proof}

The removal of the initial inradius loss by repetition is important:
a single-horizon application would leave an additive $\log(1/a)$.
Theorem~\ref{thm:rate44} is a comparison, not a duality asserting equality
or a characterization of arbitrary hidden width. The Hankel normal
form in Proposition~\ref{prop:hankel-form} permits general future-response
polytopes, whereas \eqref{eq:enclosure-profile} uses one common polytope
in the observed vector space.

\begin{corollary}[Nonuniform enclosure chains]\label{cor:chain44}
Suppose $P_t\subset\Ball^D$ has at most $K_t$ generators,
$a\Ball^D\subset P_0$, and $U_bP_{t-1}\subset\lambda_tP_t$ for every
command. Then
\[
 \mathcal I_N(K_\bullet)\le\log(1/a)+\sum_{t=1}^N\log\lambda_t.
\]
\end{corollary}
\begin{proof}
Use the cut-dependent convex rows and the initialization with mean $au$.
The same induction gives mean $a(\prod_t\lambda_t)^{-1}U_wu$.
Choose the independent packet law attaining the interval recurrence.
The pairing argument bounds the final amplitude below by this factor,
while the packet theorem bounds it above by
$\exp(-\mathcal I_N(K_\bullet))$. Taking logarithms proves the claim.
\end{proof}
